% ============================================================
% config/00_preambule.tex
% Packages et configuration générale
% ============================================================

\usepackage[french]{babel}
\usepackage{iftex}

% --- Police : pdfLaTeX (portable) vs XeLaTeX/LuaLaTeX (Times réel) ---
% LIMITE ANTICIPEE : sur Overleaf ou une machine sans les polices
% Microsoft installées, fontspec + "Times New Roman" échoue à la
% compilation. Le mode pdfLaTeX + mathptmx (ci-dessous) est le choix
% par défaut car il fonctionne partout sans dépendance externe.
\newif\ifusetruetnr
\usetruetnrfalse % passer a \usetruetnrtrue si Times New Roman est
                  % installee ET que vous compilez en XeLaTeX/LuaLaTeX

\newif\ifusexetexorluatex
\ifxetex
  \usexetexorluatextrue
\else
  \ifluatex
    \usexetexorluatextrue
  \else
    \usexetexorluatexfalse
  \fi
\fi

\ifusexetexorluatex
  \usepackage{fontspec}
  \ifusetruetnr
    \setmainfont{Times New Roman}
  \else
    \setmainfont{Liberation Serif} % clone metrique libre, tres proche
  \fi
  \newfontfamily\codefont{JetBrains Mono}[Scale=0.9]
\else
  \usepackage{mathptmx} % clone metrique de Times (Nimbus Roman)
  \usepackage[T1]{fontenc}
\fi

% --- Mise en page ---
% headheight/footskip augmentes pour laisser la place aux motifs
% ornementaux ajoutes en en-tete/pied de page (config/04_ornements.tex).
\usepackage[a4paper, margin=2.5cm, headheight=20pt, footskip=36pt]{geometry}
\linespread{1.5}
\usepackage{setspace}

% --- Packages requis par les fichiers config/ (DOIVENT etre charges
% AVANT les \input ci-dessous) ---
\usepackage{xcolor}
\usepackage{titlesec}
\usepackage{titlesec}
\usepackage{fancyhdr}
\usepackage{listings}
\usepackage{graphicx}
\usepackage{booktabs}
\usepackage{tabularx}
\usepackage{float}
\usepackage{placeins} % \FloatBarrier : empeche une figure de deriver
                       % dans le mauvais chapitre (limite classique de report)
\usepackage{caption}
\captionsetup{font=small, labelfont=bf}

% --- Habillage visuel (style ENSPY / Legrand Orange Book) ---
\usepackage{tikz}
\usetikzlibrary{calc}
\usetikzlibrary{chains} % necessaire pour la bordure de page de garde (config/04_ornements.tex)
\usepackage{pgfornament} % motifs ornementaux vectoriels (page de garde, en-tetes)
\usepackage{lettrine}    % lettrine (grande premiere lettre de chapitre)

% --- Modules centralises (couleurs / ornements / styles / metadonnees) ---
% ORDRE IMPORTANT : 04_ornements avant 02_styles, qui utilise les
% commandes d'ornement dans la definition des en-tetes/pieds de page.
% ============================================================
% config/01_couleurs.tex
% Palette du rapport - POINT UNIQUE DE MODIFICATION
% Inspiree de l'identite visuelle AGT Infra (orange primary),
% adaptee a un document academique imprimable.
% Pour changer la couleur de TOUS les titres/liens d'un coup :
% ne modifier QUE reportprimary ci-dessous.
% NOTE : le package xcolor est charge dans config/00_preambule.tex,
% avant cet \input : ne pas le recharger ici.
% ============================================================

\definecolor{reportprimary}{HTML}{C2410C}   % orange soutenu (titres, liens)
\definecolor{reportsuccess}{HTML}{059669}   % vert - a reserver aux resultats valides
\definecolor{reporterror}{HTML}{E11D48}     % rouge - a reserver aux limites/problemes
\definecolor{reportwarning}{HTML}{D97706}   % ambre - a reserver aux points d'attention

\definecolor{reportgraytext}{HTML}{334155}
\definecolor{reportgraylight}{HTML}{F1F5F9}
\definecolor{reportbordergray}{HTML}{CBD5E1}
% ============================================================
% config/04_ornements.tex
% Motifs ornementaux (pgfornament) - POINT UNIQUE DE MODIFICATION
% Identifiants issus du catalogue du package pgfornament (voir
% "texdoc pgfornament" pour la planche complete des motifs disponibles).
% Changer un identifiant ici suffit a changer le motif partout ou il
% est utilise (page de garde, en-tetes, pieds de page).
% ============================================================

% --- Motifs de la bordure de page de garde ---
\newcommand{\agtOrnCorner}{24}  % motif de coin
\newcommand{\agtOrnSide}{19}    % motif de bordure repete (haut/bas)
\newcommand{\agtOrnSize}{10mm}  % taille d'un motif de bordure

% --- Motifs legers d'en-tete / pied de page (chaque page du corps) ---
\newcommand{\agtOrnHeader}{80}
\newcommand{\agtOrnFooter}{83}
\newcommand{\agtOrnHeaderSize}{5mm}
\newcommand{\agtOrnFooterSize}{4.5mm}

% --- Motif de la page de dedicace (reutilise le motif de coin,
% \agtOrnCorner, deja valide a la compilation sur la page de garde,
% plutot qu'un nouvel identifiant non teste) ---
\newcommand{\agtOrnDedicaceSize}{16mm}

% --- Bordure complete (page de garde uniquement) ---
% CORRECTION : l'ancienne version (boucle \foreach avec rotation
% manuelle des noeuds) faussait le calcul des ancres et produisait des
% blocs pleins noirs au lieu du motif repete. Remplacee par la technique
% eprouvee du template ENSPY (Yuhala, depot latex-thesis) : une chaine
% de noeuds TikZ (bibliotheque "chains") qui tourne aux quatre coins,
% bien plus fiable que des calculs de position manuels.
% LIMITE ANTICIPEE : le nombre de motifs par cote (17 en haut/bas, 25
% sur les cotes) est calibre pour une page A4 avec un motif de 10mm ;
% si la taille du motif (\agtOrnSize) change, ces nombres doivent etre
% ajustes en consequence pour ne pas deborder ou laisser un espace vide.
\newcommand{\pageborderornament}{%
  \begin{tikzpicture}[remember picture, overlay, start chain, node distance=-2mm, color=reportprimary]%
    \node (agtnw) [shift={(5mm,-5mm)}, anchor=north west, on chain] at (current page.north west) {\pgfornament[width=\agtOrnSize]{\agtOrnCorner}};%
    \foreach \i in {1,...,16} {\node [on chain] {\pgfornament[width=\agtOrnSize]{\agtOrnSide}};}%
    \node (agtne) [on chain] {\pgfornament[width=\agtOrnSize]{\agtOrnCorner}};%
    \foreach \i in {1,...,24} {\node [continue chain=going below, on chain] {\pgfornament[width=\agtOrnSize]{\agtOrnSide}};}%
    \node (agtse) [on chain] {\pgfornament[width=\agtOrnSize]{\agtOrnCorner}};%
    \foreach \i in {1,...,16} {\node [continue chain=going left, on chain] {\pgfornament[width=\agtOrnSize]{\agtOrnSide}};}%
    \node (agtsw) [on chain] {\pgfornament[width=\agtOrnSize]{\agtOrnCorner}};%
    \foreach \i in {1,...,24} {\node [continue chain=going above, on chain] {\pgfornament[width=\agtOrnSize]{\agtOrnSide}};}%
  \end{tikzpicture}%
}

% --- Lettrine coloree (raccourci) ---
% Usage dans un chapitre : \agtlettrine{P}{our mener a bien...}
\setcounter{DefaultLines}{3}
\newcommand{\agtlettrine}[2]{\lettrine[lines=3]{\color{reportprimary}#1}{#2}}
% ============================================================
% config/02_styles.tex
% Styles des titres, en-tetes/pieds de page, blocs de code
% NOTE : titlesec, fancyhdr, listings, tikz et pgfornament sont
% charges dans config/00_preambule.tex, avant cet \input.
% ============================================================

% --- Bloc de chapitre (style Legrand Orange Book) ---
% CORRECTION : le mot "CHAPITRE" etait auparavant empile A L'INTERIEUR
% du bloc colore, au-dessus du chiffre, ce qui l'ecrasait. Il est
% desormais place A L'EXTERIEUR du bloc, en label vertical a gauche,
% cale sur la meme hauteur que le carre colore (cf. reference visuelle).
\newcommand{\ChapterNumberBlock}{%
  \begin{tikzpicture}[remember picture, baseline]
    % Rectangle vertical
    \node[fill=reportprimary, minimum width=2.0cm, minimum height=2.5cm, inner sep=0pt] (numbox) {};
    
    % Numéro étiré dynamiquement à 2.4cm de hauteur (remplit le bloc)
    \node[text=white, inner sep=0pt] at (numbox.center) {%
      \resizebox{!}{2.0cm}{\usefont{OT1}{pbd}{m}{n}\bfseries\thechapter}%
    };

    % Mot "CHAPITRE" étiré à la hauteur exacte du rectangle (3.0cm)
    \node[anchor=south, rotate=90, text=reportprimary, inner sep=0pt]
      at ([xshift=-4mm]numbox.west) {%
        \resizebox{2.5cm}{!}{\rmfamily\bfseries C\,H\,A\,P\,I\,T\,R\,E}%
      };
  \end{tikzpicture}%
}

\titleformat{\chapter}[display]
  {\normalfont\filleft} % Alignement de tout le bloc et du texte à droite
  {\ChapterNumberBlock} % Bloc de numéro
  {12pt}                % Espace vertical (saut de ligne) entre le bloc et le titre
  {\Huge\bfseries\color{reportprimary}} % Style du texte du titre
\titlespacing*{\chapter}{0pt}{0pt}{30pt}

\titleformat{\section}
  {\normalfont\Large\bfseries\color{reportprimary}}
  {\thesection}{1em}{}
\titleformat{\subsection}
  {\normalfont\large\bfseries}
  {\thesubsection}{1em}{}

% --- En-tete / pied de page, avec motifs ornementaux legers ---
\pagestyle{fancy}
\fancyhf{}
\fancyhead[L]{{\color{reportprimary}\pgfornament[width=\agtOrnHeaderSize]{\agtOrnHeader}}~\nouppercase{\leftmark}}
\fancyhead[R]{\rapporttitrecourt~{\color{reportprimary}\pgfornament[width=\agtOrnHeaderSize]{\agtOrnHeader}}}
\fancyfoot[C]{{\color{reportprimary}\pgfornament[width=\agtOrnFooterSize]{\agtOrnFooter}}~\thepage~{\color{reportprimary}\reflectbox{\pgfornament[width=\agtOrnFooterSize]{\agtOrnFooter}}}}
\renewcommand{\headrulewidth}{0.4pt}

% --- Mini-sommaire de chapitre (style ENSPY) ---
% Usage : appeler juste apres \chapter{...}\label{...}
\newcommand{\minisommaire}{%
  \etocsetnexttocdepth{subsection}%
  \etocsettocstyle{\section*{Contenu}}{}%
  \localtableofcontents%
  \bigskip%
}

% --- Blocs de code ---
% LIMITE ANTICIPEE : en pdfLaTeX, JetBrains Mono n'est pas disponible ;
% \ttfamily (police a chasse fixe par defaut) sert de repli fidele a
% l'esprit, sans dependance externe. En XeLaTeX/LuaLaTeX, \codefont
% (defini en 00_preambule) utilise la vraie police si installee.
\lstdefinestyle{codeagt}{
  backgroundcolor=\color{reportgraylight},
  frame=single,
  rulecolor=\color{reportbordergray},
  basicstyle=\ttfamily\small,
  breaklines=true,
  breakatwhitespace=false,
  showstringspaces=false,
  numbers=left,
  numberstyle=\tiny\color{reportgraytext},
  captionpos=b,
  tabsize=2
}
\lstset{style=codeagt}
% ============================================================
% config/03_metadonnees.tex
% Metadonnees - POINT UNIQUE DE MODIFICATION
% Toute information reutilisee (page de garde, en-tetes...)
% doit etre definie ICI, jamais recopiee en dur ailleurs.
% ============================================================

% Titre complet : utilise uniquement sur la page de garde.
\newcommand{\rapporttitre}{\MakeUppercase{Conception et implémentation d'AGT Infra : plateforme interne de gestion d'infrastructure et de supervision}}

% Titre court : utilise dans l'en-tete de page (evite tout debordement
% lie a la longueur du titre complet ci-dessus).
\newcommand{\rapporttitrecourt}{AGT Infra, rapport de stage}

% --- Identite institutionnelle bilingue (page de garde) ---
% ATTENTION : le departement ci-dessous est repris de l'image de
% reference fournie par Josue (autre rapport ENSPY), confirme par lui
% au fil de la session. A corriger ici si ce n'est pas exact.
\newcommand{\rapportuniversitefr}{UNIVERSITÉ DE YAOUNDÉ I}
\newcommand{\rapportuniversiteen}{UNIVERSITY OF YAOUNDÉ I}
\newcommand{\rapportecolefr}{ÉCOLE NATIONALE SUPÉRIEURE POLYTECHNIQUE}
\newcommand{\rapportecoleen}{NATIONAL ADVANCED SCHOOL OF ENGINEERING}
\newcommand{\rapportdepartementfr}{DÉPARTEMENT DE GÉNIE INFORMATIQUE}
\newcommand{\rapportdepartementen}{DEPARTMENT OF COMPUTER ENGINEERING}

% --- Type de stage, auteur, entreprise, tuteur ---
\newcommand{\rapporttypestage}{PRÉ-INGÉNIEUR}
\newcommand{\rapportauteur}{MBAH MOFFO Josué}
\newcommand{\rapportentreprise}{AG Technologies}
\newcommand{\rapporttuteurentreprise}{Ing. NOMO BODIANGA Gabriel}
\newcommand{\rapporttuteuracademique}{[Nom du tuteur académique]}
\newcommand{\rapportperiode}{Du 06 juillet 2026 au 05 septembre 2026}
\newcommand{\rapportannee}{2025 - 2026}

% --- Logos (page de garde) ---
\newcommand{\rapportlogoecole}{assets/logos/logo-ENSPY.png}
\newcommand{\rapportlogoentreprise}{assets/logos/logo-AGT.png}

% --- Bibliographie ---
\usepackage[backend=biber, style=ieee, sorting=none]{biblatex}
\addbibresource{biblio.bib}

% --- Glossaire / acronymes ---
% LIMITE ANTICIPEE : necessite une passe "makeglossaries" en plus de
% pdflatex/biber. latexmk gere cela automatiquement (voir main.tex).
\usepackage[acronym, toc, nonumberlist]{glossaries}
\makeglossaries

% --- Liens cliquables ---
% Bascule impression : les liens colores genent souvent a l'impression.
\newif\ifmodeimpression
\modeimpressionfalse % passer a true pour generer une version "a imprimer"

\usepackage{hyperref}
\ifmodeimpression
  \hypersetup{colorlinks=true, linkcolor=black, citecolor=black, urlcolor=black}
\else
  \hypersetup{colorlinks=true, linkcolor=reportprimary, citecolor=reportprimary, urlcolor=reportprimary}
\fi

% --- Mini-sommaire par chapitre (style ENSPY) ---
% LIMITE ANTICIPEE : etoc doit etre charge apres hyperref pour que les
% liens du mini-sommaire fonctionnent correctement.
\usepackage{etoc}

% --- Veuves / orphelines ---
\clubpenalty=10000
\widowpenalty=10000
\displaywidowpenalty=10000